\chapter{کدگذاری موقعیت}

\section{جایگاه این بحث در درس}
در جلسه‌های پیشین بحث \eng{Embedding} را به پایان رساندیم و آخرین موضوع، معماری \eng{Sequence-to-Sequence} برای پردازش نویسه‌ها (مسئلهٔ \eng{G2P}) بود. پیش از ادامهٔ آن معماری، باید یکی از مشکلات مهم \eng{Embedding} پویا (بافتی) و راه‌حل آن را بشناسیم: \term{کدگذاری موقعیت}{Positional Encoding}. پس از آن به معماری قبلی برمی‌گردیم و اجزای بیشتری به شبکه اضافه می‌کنیم تا معماری توانمندتری داشته باشیم. مفاهیم \eng{Query}، \eng{Key}، \eng{Value} و \eng{Self-Attention} را در فصل قبل (بخش «راه‌های ساخت \eng{Embedding} بافتی» و بخش \eng{Cross-Attention}) دیدیم و اینجا فقط به اندازهٔ نیاز مرور می‌کنیم.

\begin{exambox}{نکتهٔ روی تخته}
کدگذاری موقعیت \emph{فقط در \eng{Embedding} پویا معنا پیدا می‌کند}؛ \eng{Embedding} ایستا مستقل از بافت (\eng{context-free}) است و در آن اصلاً مسئلهٔ «نقش کلمه نسبت به کلمه‌های دیگر» مطرح نیست.
\end{exambox}

\section{مسئلهٔ ترتیب در \eng{Attention}}

\subsection{دو جمله با کلمات یکسان}
دو جملهٔ زیر را در نظر بگیرید:
\begin{itemize}
  \item «ایران آمریکا را شکست داد.»
  \item «آمریکا ایران را شکست داد.»
\end{itemize}
کلمه‌ها یکسان‌اند ولی معنا کاملاً فرق می‌کند: در جملهٔ اول ایران برنده است و در جملهٔ دوم آمریکا. فرض کنید کانون توجه کلمهٔ «شکست» باشد. در \eng{Embedding} پویا انتظار داریم این کلمه به بقیهٔ جمله «نگاه کند»، اطلاعات لازم را بگیرد و نمایشی وابسته به بافت بسازد. چون نقش کلمه‌ها در دو جمله متفاوت است، نمایش «شکست» در این دو جمله باید متفاوت باشد. در \eng{Embedding} ایستا، کلمهٔ «شکست» در ده جملهٔ مختلف اساساً همان بردار ثابت را دارد و این تمایز اصلاً ساخته نمی‌شود؛ ولی در \eng{Embedding} پویا همین تمایز هدف اصلی است.

\subsection{مرور \eng{Query}، \eng{Key}، \eng{Value}}
در رویکرد موازی (\eng{Self-Attention}) کلمهٔ هدف («شکست») نقش \eng{Query} را دارد و کلمه‌های دیگر هرکدام یک \eng{Key} و یک \eng{Value} دارند. مدل با مقایسهٔ \eng{Query} با \eng{Key}ها وزن توجه $\alpha_i$ را برای هر کلمه می‌سازد و بردار بافت را به‌صورت ترکیب وزن‌دار \eng{Value}ها به دست می‌آورد:
\[
\text{\lr{Context Vector}} \;=\; \sum_i \alpha_i V_i .
\]

\begin{figure}[htbp]
\centering
\resizebox{0.97\textwidth}{!}{%
\begin{tikzpicture}[rtl, 
  w/.style={font=\Large,inner sep=3pt},
  vec/.style={font=\large,text=mainblue},
  hl/.style={draw=red!70!black,thick,ellipse,inner sep=3pt}]
  % sentence 1 (right to left)
  \node[w] at (13.0,0) {\RL{جملهٔ ۱}};
  \node[w] (a1) at (10.2,0) {\RL{ایران}};
  \node[w] (a2) at (7.8,0) {\RL{آمریکا}};
  \node[w] at (5.6,0) {\RL{را}};
  \node[w,hl] (a4) at (3.4,0) {\RL{شکست}};
  \node[w] at (1.0,0) {\RL{داد}};
  \node[vec] at (10.2,-1.3) {$k_1=\binom{2}{0}$};
  \node[vec] at (7.8,-1.3) {$k_2=\binom{0}{2}$};
  \node[vec] at (3.4,1.4) {$Q=\binom{1}{1}$};
  \draw[-{Stealth},mainblue] (3.4,0.9)--(3.4,0.55);
  % sentence 2
  \node[w] at (13.0,-4.6) {\RL{جملهٔ ۲}};
  \node[w] (b1) at (10.2,-4.6) {\RL{آمریکا}};
  \node[w] (b2) at (7.8,-4.6) {\RL{ایران}};
  \node[w] at (5.6,-4.6) {\RL{را}};
  \node[w,hl] (b4) at (3.4,-4.6) {\RL{شکست}};
  \node[w] at (1.0,-4.6) {\RL{داد}};
  \node[vec] at (10.2,-5.9) {$k_2=\binom{0}{2}$};
  \node[vec] at (7.8,-5.9) {$k_1=\binom{2}{0}$};
  \node[vec] at (3.4,-3.2) {$Q=\binom{1}{1}$};
  \draw[-{Stealth},mainblue] (3.4,-3.7)--(3.4,-4.05);
  % scores
  \node[vec,align=center] at (3.4,-1.5) {$Q\!\cdot\! k_1=2,\ Q\!\cdot\! k_2=2$\\[2pt]$\alpha_1=\alpha_2=0{,}5$};
  \node[vec,align=center] at (3.4,-6.1) {$Q\!\cdot\! k_1=2,\ Q\!\cdot\! k_2=2$\\[2pt]$\alpha_1=\alpha_2=0{,}5$};
  \draw[gray!60,dashed] (-0.8,-2.7)--(14.6,-2.7);
\end{tikzpicture}}
\caption{مسئلهٔ ترتیب: کلمهٔ «شکست» (دایره‌شده) در هر دو جمله همان $Q=(1,1)$ را دارد و \eng{Key} کلمه‌های دیگر هم یکسان است؛ فقط جای آن‌ها عوض شده است. در هر دو جمله $Q\cdot k_1=Q\cdot k_2=2$ و وزن‌ها برابر می‌شوند، پس بدون اطلاعات موقعیت مدل تفاوتی نمی‌بیند. (بازسازی‌شده از تختهٔ کلاس)}
\end{figure}

\subsection{مثال عددی ساده‌شدهٔ کلاس}
روی تخته، \eng{Query} کلمهٔ «شکست» را بردار $q=(1,1)$ و \eng{Key} دو کلمهٔ دیگر را (طبق خوانش تخته) $k_{1}=(2,0)$ و $k_{2}=(0,2)$ گرفتیم. هر دو ضرب داخلی برابر ۲ می‌شوند:
\[
q\cdot k_1 = 1\cdot 2+1\cdot 0 = 2,\qquad q\cdot k_2 = 1\cdot 0+1\cdot 2 = 2 ,
\]
پس امتیازها برابرند و وزن‌ها $\alpha_1=\alpha_2=0{,}5$ می‌شوند. اگر \eng{Value} هر دو را ۱ فرض کنیم، بردار بافت $C = 0{,}5\times 1 + 0{,}5\times 1 = 1$ است. حالا جای دو کلمه را عوض کنید: \eng{Key} و \eng{Value} هر کلمه عوض نشده و فقط ترتیب فهرست جابه‌جا شده است؛ جمع وزن‌دار دقیقاً همان ۱ را می‌دهد. یعنی مدل نمی‌تواند بفهمد کدام کلمه زودتر آمده است.

پرسش: در دو جمله چه چیزی تغییر کرده است؟ پاسخ: \emph{موقعیت} کلمه‌ها. تا اینجا اطلاعات موقعیت را وارد \eng{Embedding} نکرده‌ایم، پس در رویکرد موازی مدل فقط «مجموعه‌ای» از نمایش کلمه‌ها را می‌بیند و ترتیب را صریحاً نمی‌داند. راه‌حل، افزودن اطلاعات موقعیت به نمایش هر کلمه است.

\begin{definitionbox}{کدگذاری موقعیت (\eng{Positional Encoding})}
افزودن اطلاعات «این عنصر چندمین عنصر دنباله است» به نمایش هر عنصر، تا مدل علاوه بر اینکه بداند هر کلمه \emph{چیست}، بداند \emph{کجا} قرار گرفته است. مثلاً در «ایران آمریکا را شکست داد»: «ایران» موقعیت ۱، «آمریکا» موقعیت ۲، «را» موقعیت ۳ و «شکست» موقعیت ۴. با جابه‌جایی کلمه‌ها، اطلاعات موقعیت هم باید متناسب تغییر کند.
\end{definitionbox}

\begin{correctionbox}{دقت: «بی‌توجه به ترتیب» یعنی چه؟}
در مثال کلاس برای ساده‌سازی گفته شد که ترتیب کلمه‌ها «اثری ندارد». دقیق‌تر: \eng{Self-Attention} بدون کدگذاری موقعیت \term{هم‌متغیر نسبت به جایگشت}{Permutation-Equivariant} است؛ یعنی اگر ورودی را جابه‌جا کنیم، خروجی هر عنصر تغییر نمی‌کند و فقط سطرها همان جابه‌جایی را می‌بینند. بردار بافت یک \eng{Query} مشخص نسبت به جابه‌جایی فهرست \eng{Key}/\eng{Value}ها کاملاً ثابت است (همان چیزی که در مثال بالا دیدیم). این با «ترتیب بی‌اثر است» به یک معنا یکی است، ولی نباید با «\eng{Attention} هیچ‌چیز از بافت نمی‌فهمد» اشتباه گرفته شود: کلمه‌ها هنوز محتوای هم را می‌بینند و تنها اطلاعات \emph{ترتیب} از دست می‌رود.
\end{correctionbox}

\begin{addedbox}{نمایش عددی: بدون و با کدگذاری موقعیت}
سه بردار دوبعدی داریم: $A=(1,0)$، $B=(0,1)$ و بردار کلمهٔ هدف $S=(1,1)$. جملهٔ اول به ترتیب $A,B,S$ و جملهٔ دوم به ترتیب $B,A,S$ است. برای سادگی $W_Q=W_K=W_V=I$ و امتیازها بر $\sqrt{2}$ تقسیم می‌شوند. خروجی \eng{Attention} برای $S$ (که همهٔ سه عنصر را می‌بیند):
\begin{center}
\begin{tabular}{lcc}
\toprule
 & ترتیب $A,B,S$ & ترتیب $B,A,S$ \\
\midrule
بدون کدگذاری موقعیت & $(0{,}752;\,0{,}752)$ & $(0{,}752;\,0{,}752)$ \\
با افزودن $\mathrm{PE}$ سینوسی ($d=2$) & $(1{,}502;\,0{,}867)$ & $(1{,}751;\,0{,}660)$ \\
\bottomrule
\end{tabular}
\end{center}
بدون موقعیت دو خروجی یکی است؛ با افزودن $\mathrm{PE}$ (ردیف‌های ۰ تا ۲ ماتریس سینوسی، بخش‌های بعد) دو جمله از هم متمایز می‌شوند.
\end{addedbox}

\section{آیا مدل‌های بازگشتی هم به کدگذاری موقعیت نیاز دارند؟}
\label{sec:ch10-rnn}
\begin{assignmentbox}{بحث کنید (روی تخته)}
بحث کنید که آیا در روش بازگشتی هم به پوزیشن انکدینگ نیاز داریم؟
\end{assignmentbox}

پاسخ مطرح‌شده در کلاس: در رویکرد بازگشتی (\eng{RNN}/\eng{LSTM}) عنصرهای دنباله \emph{به‌ترتیب} پردازش می‌شوند. مدل کلمهٔ اول را می‌بیند و یک حالت پنهان می‌سازد؛ با رسیدن کلمهٔ دوم، حالت پنهان اطلاعات هر دو را در خود دارد؛ و همین‌طور تا انتها. پس ترتیب در خودِ فرایند محاسبه نهفته است. مثلاً اگر در یک دنبالهٔ نویسه‌ای حرف «ت» دو بار تکرار شود، حالت پنهان هنگام دیدن «ت» دوم با حالت پنهان هنگام دیدن «ت» اول یکسان نیست، چون اولی چیزی از گذشته ندارد و دومی اطلاعات نویسه‌های قبلی (از جمله همان «ت» اول) را حمل می‌کند. بنابراین در رویکرد بازگشتیِ بررسی‌شده در کلاس نیازی به افزودن صریح کدگذاری موقعیت، به شکلی که برای رویکرد موازی لازم است، نداریم. در مقابل، در رویکرد موازی همهٔ عنصرها هم‌زمان پردازش می‌شوند و ترتیب خودبه‌خود وارد محاسبهٔ \eng{Attention} نمی‌شود.

\begin{table}[htbp]
\centering
\caption{مقایسهٔ منشأ اطلاعات ترتیب در دو رویکرد ساخت \eng{Embedding} بافتی.}
\begin{tabular}{lcc}
\toprule
 & ترتیبی (\eng{RNN}/\eng{LSTM}) & موازی (\eng{Self-Attention}) \\
\midrule
نحوهٔ پردازش & گام‌به‌گام & همهٔ عنصرها هم‌زمان \\
منشأ اطلاعات ترتیب & خودِ بازگشت (حالت پنهان) & باید صریحاً افزوده شود \\
نیاز به کدگذاری موقعیت & خیر (در بحث کلاس) & بله \\
محدودیت اصلی & فراموشی در دنباله‌های بلند، پردازش غیرموازی & نیاز به تعریف موقعیت \\
\bottomrule
\end{tabular}
\end{table}

\begin{addedbox}{ظرایف این پرسش}
\begin{itemize}
  \item اطلاعات ترتیب در \eng{RNN} «ضمنی» است: حالت پنهان نشان نمی‌دهد «الان در موقعیت ۱۷ هستیم»، بلکه تاریخچه‌ای فشرده از آنچه دیده شده است. برای بسیاری از کارها همین کافی است. محدودیت دیگری هم وجود دارد (فراموشی اطلاعات قدیمی در دنباله‌های بلند) که جدا بررسی می‌شود.
  \item اگر روی خروجی‌های \eng{RNN} یک لایهٔ \eng{Attention} بگذاریم (مانند معماری‌های رمزگذار--رمزگشا با \eng{Attention} پیش از \eng{Transformer})، حالت‌های رمزگذار پیشاپیش اطلاعات ترتیب را حمل می‌کنند و باز هم کدگذاری موقعیت صریح لازم نیست.
  \item در تصویر تخته برای مثال عددی (بخش \ref{sec:ch10-example}) در معماری رمزگذار--رمزگشا، هم‌زمان با \eng{Embedding} نویسه‌ها (\eng{CE}) یک \eng{PE} نوشته شده است؛ یعنی در معماری \eng{Attention}‌محور، کدگذاری موقعیت باید حضور داشته باشد.
\end{itemize}
\end{addedbox}

\section{سه پرسش طراحی}
در کلاس مسئله به سه پرسش شکسته شد که روی تخته به‌صورت یک نمودار دیده می‌شود:
\begin{enumerate}
  \item \textbf{محل}: کدگذاری موقعیت را در کدام بخش معماری قرار دهیم؟
  \item \textbf{نحوهٔ اعمال}: اطلاعات موقعیت را چگونه به \eng{Embedding} متصل کنیم؟ (جمع یا کانکت)
  \item \textbf{نحوهٔ مقداردهی}: مقدارهای کدگذاری موقعیت را چگونه تعیین کنیم؟ (یادگیری یا فرمول)
\end{enumerate}

\begin{figure}[htbp]
\centering
\resizebox{\textwidth}{!}{%
\begin{tikzpicture}[rtl, 
  box/.style={draw,rounded corners=4pt,fill=softgreen,minimum height=1.5cm,align=center,font=\normalsize,text width=4.2cm},
  note/.style={draw,rounded corners=3pt,fill=boxgray,align=center,font=\small,text width=6.2cm,inner sep=5pt},
  ar/.style={-{Stealth[length=2.5mm]},thick}]
  \node[draw,diamond,aspect=2.4,fill=softgreen,font=\large,inner sep=2pt,align=center] (root) at (0,0) {\lr{positional}\\\lr{encoding}};
  \node[box,draw=orange,very thick] (loc) at (6.5,-2.8) {\RL{محل پوزیشنال انکدینگ در معماری؟}};
  \node[box] (app) at (0,-4.2) {\RL{نحوه اعمال؟}};
  \node[box] (val) at (-6.5,-2.8) {\RL{نحوه مقداردهی؟}};
  \draw[ar] (root.east)--(loc.north);
  \draw[ar] (root.south)--(app.north);
  \draw[ar] (root.west)--(val.north);
  % apply options
  \node[draw,rounded corners=3pt,fill=softyellow,font=\normalsize,minimum width=2.2cm] (sum) at (1.5,-6.6) {\RL{جمع}};
  \node[draw,rounded corners=3pt,fill=softyellow,font=\normalsize,minimum width=2.2cm] (cat) at (-1.5,-6.6) {\RL{کانکت}};
  \draw[ar] (app.south)--(sum.north); \draw[ar] (app.south)--(cat.north);
  % value options (stacked)
  \node[note] (learn) at (-6.5,-5.9) {\lr{Learnable Position Embedding Matrix}\\ \RL{پارامتر قابل‌آموزش و گرادیان باعث به‌روزرسانی آن‌ها می‌شود. اکتیویشن نداریم، ضرب ماتریسی نداریم (با لایهٔ خطی اشتباه نشود)}};
  \node[note] (form) at (-6.5,-9.9) {\RL{فرمول تولید عددی با مین و ماکس مشخص بر اساس شمارهٔ موقعیت و شمارهٔ بعد}\\[3pt] $PE(pos,j)=\sin\big(g(pos,j)\big)$\\[3pt]{\footnotesize \RL{\lr{sin(?)}: هیچ دو خانه‌ای زاویه یکسان ندارند}}};
  \draw[ar] (val.south)--(learn.north);
  \draw[ar] (val.west) -- ++(-1.0,0) |- (form.west);
\end{tikzpicture}}
\caption{نمودار تخته: سه پرسش طراحی کدگذاری موقعیت (در این بازسازی به ترتیب خواندن فارسی از راست به چپ چیده شده‌اند). پرسش محل در معماری روی تخته دایره شده بود؛ برای مقداردهی دو راه (ماتریس یادگرفتنی و فرمول) یادداشت شده بود. (بازسازی‌شده از تختهٔ کلاس)}
\end{figure}

\subsection{محل قرارگیری در معماری}
فرض کنید ابتدا یک لایهٔ \eng{Embedding} داریم که برای هر توکن یک بردار اولیه می‌سازد و سپس می‌خواهیم با \eng{Attention} این بردار را با توجه به بقیهٔ عنصرها بهبود دهیم. اطلاعات موقعیت باید \emph{بعد از \eng{Embedding} اولیه و پیش از ورود به محاسبات \eng{Attention}} اضافه شود، چون هنگام مقایسهٔ \eng{Query} و \eng{Key} مدل باید بتواند تفاوت جایگاه عنصرها را هم ببیند.

\begin{figure}[htbp]
\centering
\begin{tikzpicture}[rtl, 
  blk/.style={draw,rounded corners=3pt,minimum width=5.2cm,minimum height=0.85cm,align=center,font=\small},
  ar/.style={-{Stealth[length=2mm]},thick}]
  \node[blk,fill=gray!12,minimum width=7.6cm] (tok) at (0,0) {\RL{توکن‌های ورودی (کلمه، نویسه، واج، \lr{byte})}};
  \node[blk,fill=blue!10] (emb) at (0,1.6) {\lr{Embedding} اولیه: $E\in\mathbb{R}^{T\times d}$};
  \node[blk,fill=green!15,minimum width=2.8cm] (pe) at (4.6,3.2) {\RL{$P$: کدگذاری موقعیت}};
  \node[draw,circle,fill=white,inner sep=1pt] (plus) at (0,3.2) {$+$};
  \node[blk,fill=orange!15,minimum width=6.4cm] (att) at (0,4.8) {\lr{Attention}: $Q$، $K$، $V$ از $X=E+P$};
  \node[blk,fill=blue!10] (ctx) at (0,6.4) {\RL{\lr{Embedding} بافتی (پویا)}};
  \draw[ar] (tok)--(emb); \draw[ar] (emb)--(plus); \draw[ar] (plus)--(att); \draw[ar] (att)--(ctx);
  \draw[ar] (pe)--(plus);
\end{tikzpicture}
\caption{جای کدگذاری موقعیت در معماری: پس از \eng{Embedding} اولیه و پیش از \eng{Attention}. (شکل افزوده)}
\end{figure}

\begin{figure}[htbp]
\centering
\resizebox{0.8\textwidth}{!}{%
\begin{tikzpicture}[rtl, ar/.style={-{Stealth[length=2.5mm]},thick,mainblue}]
  % E column
  \draw[fill=softblue] (0,0) rectangle (0.8,6);
  \node[above] at (0.4,6.1) {\Large $E$};
  \node[left=3pt] at (0,3) {\Large 512};
  \node[below] at (0.4,-0.1) {\lr{Embedding}};
  % P matrix
  \draw[fill=white] (3.5,0) rectangle (11.5,6);
  \foreach \x/\l in {3.5/1, 4.1/2}{
    \fill[softgreen] (\x,0) rectangle (\x+0.6,6);
    \draw (\x,0) rectangle (\x+0.6,6);
    \node[above] at (\x+0.3,6.1) {\large \l};
  }
  \fill[green!35] (7.7,0) rectangle (8.3,6);
  \draw (7.7,0) rectangle (8.3,6);
  \node[above] at (8.0,6.1) {\large $p$};
  \node[above] at (6.0,6.1) {\large $\cdots$};
  \node[above] at (10.0,6.1) {\large $\cdots$};
  \fill[softgreen] (10.9,0) rectangle (11.5,6);
  \draw (10.9,0) rectangle (11.5,6);
  \node[above] at (11.2,6.1) {\large 100};
  \node[left=3pt] at (3.5,3) {\Large 512};
  \node[below] at (7.5,-0.1) {\Large $P$};
  \node[below,font=\small] at (7.5,-1.0) {\RL{۱۰۰ ستون؛ هر ستون یک موقعیت}};
  % arrow
  \draw[ar] (8.0,6.9) to[out=100,in=80] node[midway,above]{\Large $+$} (0.4,6.9);
  \node[font=\small] at (4.3,8.6) {\RL{ستون موقعیت $p$}};
\end{tikzpicture}}
\caption{طرح تخته برای جمع $E+P$: بردار \eng{Embedding} ۵۱۲ بعدی (ستون سمت چپ) با ستونِ موقعیت مورد نظر از ماتریس موقعیت (۵۱۲ سطر و ۱۰۰ ستون) جمع می‌شود. روی تخته ماتریس به‌صورت ترانهادهٔ $P\in\mathbb{R}^{100\times512}$ رسم شده بود. (بازسازی‌شده از تختهٔ کلاس)}
\end{figure}

\begin{exambox}{انعطاف معماری}
کدگذاری موقعیت یک \emph{ماژول} است و محل و شیوهٔ استفاده از آن به معماری وابسته است. در طراحی شبکه باید بدانیم در کدام مرحله به اطلاعات موقعیت نیاز داریم، به‌ویژه اگر \eng{Attention} بر اساس نمایش‌های ورودی محاسبه می‌شود. استاد اشاره کرد ممکن است در امتحان سؤال طراحی معماری مطرح شود، مثلاً اینکه چگونه اطلاعات موقعیت را در بخش مناسبی از یک شبکه وارد کنیم.
\end{exambox}

\subsection{نحوهٔ اعمال: کانکت یا جمع}
\paragraph{کانکت (\eng{Concatenation}).} بردار موقعیت جداگانه‌ای (مثلاً ۱۰ بعدی) می‌سازیم و به انتهای \eng{Embedding} ۵۱۲ بعدی می‌چسبانیم: $512+10=522$. مشکل: ابعاد نمایش بزرگ‌تر می‌شود، در حالی که انتخاب ابعاد \eng{Embedding} خود یک تصمیم مهم طراحی است و نمی‌خواهیم با هر اطلاعات جدید ابعاد را افزایش دهیم.

\paragraph{جمع (\eng{Addition}).} بردار موقعیت را با \eng{Embedding} جمع می‌کنیم؛ پس باید دقیقاً هم‌اندازهٔ آن باشد (برای \eng{Embedding} ۵۱۲ بعدی، بردار موقعیت هم ۵۱۲ بعدی):
\[
\boxed{X \;=\; E + P}
\]
نمایش نهایی همچنان ۵۱۲ بعدی می‌ماند. این روشی است که در کلاس برای ادامهٔ بحث انتخاب شد.

\begin{table}[htbp]
\centering
\caption{مقایسهٔ کانکت و جمع برای ترکیب $E$ (با $d=512$) و بردار موقعیت. (جدول افزوده)}
\begin{tabular}{lcc}
\toprule
ویژگی & کانکت & جمع \\
\midrule
بُعد بردار موقعیت & $d_p$ دلخواه (مثلاً ۱۰) & باید برابر $d$ باشد (۵۱۲) \\
بُعد نمایش نهایی & $d+d_p=522$ & $d=512$ \\
تأثیر روی لایه‌های بعد & ابعاد ورودی بزرگ‌تر و پارامتر بیشتر & بدون تغییر \\
جداماندن محتوا و موقعیت & کاملاً در مختصات جدا & در یک فضای مشترک ترکیب می‌شوند \\
\bottomrule
\end{tabular}
\end{table}

\begin{addedbox}{آیا جمع، اطلاعات را «خراب» نمی‌کند؟}
در نگاه اول جمع شبیه آمیختن محتوا و موقعیت است. با این حال: (الف) در ابعاد زیاد، دو بردار تقریباً «در جهت‌های متفاوتی» قرار می‌گیرند و شبکه می‌تواند هر دو را از مجموع بازیابی کند؛ (ب) اگر $W$ ماتریس نخستین تبدیل خطی بعدی باشد، $W(E+P)=WE+WP$ است، و این برابر با کانکت بردارها و ضرب در ماتریسی است که دو نیمهٔ آن $W$ ‏یکسان‌اند؛ یعنی جمع حالت خاصی از کانکت با وزن‌های مشترک است و پارامتر اضافی نمی‌خواهد. این استدلال شهودی است، نه برهانی که برای هر معماری درست باشد.
\end{addedbox}

\subsection{ماتریس کدگذاری موقعیت}
فرض کنید ابعاد \eng{Embedding} برابر ۵۱۲ باشد و بخواهیم موقعیت‌ها را تا طول ۱۰۰ پوشش دهیم. ماتریس موقعیت چنین است:
\[
\boxed{P\in\mathbb{R}^{100\times 512}}
\]
هر \emph{سطر} مربوط به یک موقعیت و هر \emph{ستون} مربوط به یک بعد از \eng{Embedding} است. اگر کلمه‌ای در موقعیت دهم باشد، سطر دهم $P$ را برمی‌داریم و با \eng{Embedding} همان کلمه جمع می‌کنیم:
\[
X_{pos} \;=\; E_{token} + P_{pos}.
\]
(روی تخته ماتریس به‌صورت $512$ سطر و ۱۰۰ ستون، یعنی ترانهادهٔ این نمایش، کشیده شده بود: هر ستون یک موقعیت. این فقط قرارداد نمایش است.) تا اینجا فقط اندازهٔ ماتریس معلوم است؛ مقدارها را از کجا بیاوریم؟ دو راه مطرح شد: ماتریس را \emph{یاد بگیریم} یا با یک \emph{فرمول ثابت} بسازیم.

\section{روش اول: کدگذاری موقعیت یادگرفتنی}

\subsection{ایده}
مقدارهای $P$ را از پیش با فرمول تعیین نمی‌کنیم؛ ماتریس را مجموعه‌ای از پارامترهای قابل‌آموزش می‌گیریم (\eng{Learnable Position Embedding Matrix}) و مدل مقدارهای آن را در آموزش یاد می‌گیرد. این همان ایدهٔ \eng{Embedding} قابل‌آموزش است، در برابر \eng{One-Hot} که با یک قاعدهٔ ثابت ساخته می‌شود. به‌جای اینکه خودمان بگوییم موقعیت اول چه اعدادی دارد، این را به مدل می‌سپاریم. در مثال کلاس، $100\times 512=51{,}200$ پارامتر اضافه می‌شود.

\subsection{آیا لایهٔ خطی است؟ فعال‌ساز دارد؟}
روی تخته نوشته شد: «اکتیویشن نداریم، ضرب ماتریسی نداریم (با لایه خطی اشتباه نشود)». این ماتریس مجموعه‌ای از پارامترهاست که \emph{مستقیماً} به \eng{Embedding} اولیه اضافه می‌شود: ضرب ماتریسی، \eng{Bias} یا تابع فعال‌ساز روی آن اعمال نمی‌شود. فعال‌ساز برای ایجاد رفتار غیرخطی در لایه‌ها به کار می‌رود، ولی اینجا هدف ساخت یک نگاشت غیرخطی نیست؛ فقط می‌خواهیم برای هر موقعیت یک بردار داشته باشیم و مقدارهای آن در آموزش اصلاح شوند.

\begin{correctionbox}{دقت: جدول جست‌وجو}
در پیاده‌سازی، این ماتریس همان یک \term{جدول جست‌وجو}{Lookup Table} است، دقیقاً مثل ماتریس \eng{Embedding} کلمه‌ها: شمارهٔ موقعیت، سطر مربوط را بازمی‌گرداند. از نظر ریاضی معادل ضرب بردار \eng{One-Hot} موقعیت در $P$ است، اما نه ضرب واقعاً اجرا می‌شود و نه \eng{Bias} یا فعال‌ساز وجود دارد. پس دو عبارت «لایهٔ خطی نیست» و «در حد یک ماتریس پارامتر است» هر دو درست‌اند.
\end{correctionbox}

\subsection{آموزش از راه پس‌انتشار}
پرسش دانشجو: اگر فقط یک ماتریس پارامتر است و لایهٔ خطی هم نیست، چطور یاد گرفته می‌شود؟ پاسخ: مثل بقیهٔ پارامترهای شبکه، با \eng{Backpropagation}. خروجی جمع $X=E+P$ وارد بقیهٔ شبکه می‌شود و در انتها به یک تابع خطا (\eng{Loss}) می‌رسیم. این خطا نسبت به خروجی نهایی و هدف همان مسئله حساب می‌شود (رگرسیون با \eng{MSE}، دسته‌بندی، یا حتی چند تابع هزینهٔ هم‌زمان)، نه نسبت به خود \eng{Embedding}. گرادیان آن با قاعدهٔ زنجیره‌ای به عقب برمی‌گردد و چون $P$ در محاسبهٔ خروجی نقش داشته، به $P$ هم می‌رسد:
\[
\text{\lr{Loss}}\;\xrightarrow{\;\text{\lr{Backprop}}\;}\; E,\;P .
\]
چون $X_t=E_{token_t}+P_t$ یک جمع است، $\dfrac{\partial L}{\partial P_t}=\dfrac{\partial L}{\partial X_t}$.

\begin{figure}[htbp]
\centering
\begin{tikzpicture}[rtl, x=0.55cm,y=0.42cm]
  \foreach \r in {1,...,8}{\foreach \c in {1,...,8}{\draw[gray!50] (\c-1,-\r+1) rectangle (\c,-\r);}}
  \foreach \r in {1,2,3}{\foreach \c in {1,...,8}{\fill[green!35,opacity=.7] (\c-1,-\r+1) rectangle (\c,-\r);}}
  \foreach \c in {1,...,8}{\fill[gray!15] (\c-1,-7+1) rectangle (\c,-8);}
  \node[left] at (-0.2,-0.5) {\lr{pos 1}}; \node[left] at (-0.2,-1.5) {\lr{pos 2}}; \node[left] at (-0.2,-2.5) {\lr{pos 3}};
  \node[left] at (-0.2,-4.5) {$\vdots$}; \node[left] at (-0.2,-7.5) {\lr{pos 100}};
  \node[right,align=left,font=\small] at (8.4,-1.5) {سطرهای دیده‌شده: گرادیان غیرصفر\\و به‌روزرسانی مؤثر};
  \node[right,align=left,font=\small] at (8.4,-7.5) {سطرهای دیده‌نشده: گرادیان صفر\\(مقدار اولیه تقریباً بدون تغییر)};
\end{tikzpicture}
\caption{اگر دادهٔ آموزشی فقط توالی‌هایی تا طول ۳ داشته باشد، فقط سه سطر نخست $P$ گرادیان می‌گیرند. (شکل افزوده)}
\end{figure}

\subsection{آیا بردارهای موقعیت حتماً متمایزند؟}
در حالت یادگرفتنی شرط صریحی برای متمایز بودن سطرهای $P$ نمی‌گذاریم و به مدل نمی‌گوییم موقعیت اول و دوم حتماً فرق کنند. انتظار این است که با داده، تابع خطا و ابعاد زیاد، مدل نمایش مناسبی بیاموزد؛ استاد اشاره کرد درباره تضمین دقیق این موضوع می‌شود بیشتر بررسی کرد. (در عمل سطرها با مقدار اولیهٔ تصادفی شروع می‌شوند و بنابراین از ابتدا هم متفاوت‌اند؛ مسئله این است که سطرِ آموزش‌ندیده معنای مفیدی ندارد، نه اینکه لزوماً با سطر دیگر برابر باشد.)

\section{روش دوم: کدگذاری موقعیت ثابت (فرمولی)}

\subsection{ورودی‌های فرمول}
اگر نخواهیم $P$ را یاد بگیریم، باید قاعده‌ای برای تولید آن بنویسیم. چون هدف این است که مدل بتواند موقعیت‌های مختلف را از هم تشخیص دهد (مثلاً دو جملهٔ «شکارچی پلنگ را زد» و «پلنگ شکارچی را زد»)، نمایش موقعیت‌های مختلف باید \emph{قابل‌تمایز} باشد؛ اگر موقعیت اول و دوم نمایش یکسانی بگیرند، مشکل پابرجا می‌ماند. هر خانهٔ ماتریس به دو چیز وابسته است: \emph{شمارهٔ موقعیت} ($pos$) و \emph{شمارهٔ بعد} ($j$):
\[
P(pos,j)=f(pos,j).
\]

\subsection{فرمول سینوسی}
روی تخته نوشته شد: «فرمول تولید عددی با مین و مکس مشخص بر اساس شمارهٔ موقعیت و شمارهٔ بعد»:
\[
PE(pos,j)=\sin\big(g(pos,j)\big),
\]
که در آن $g(pos,j)$ یک زاویه است که از موقعیت و بعد ساخته می‌شود. چون $-1\le\sin(x)\le 1$ (و همین بازه برای کسینوس)، مقدارها \emph{کراندار} و عددی‌اند. استاد وارد جزئیات فرمول نشد و تأکید را روی منطق گذاشت، نه حفظ فرمول؛ فقط اشاره کرد در فرمول متداول از سینوس و کسینوس با الگوهای متفاوت برای اندیس‌های زوج و فرد استفاده می‌شود.

\boardfigure{0.75}{ch10-board-sinusoid-curves.png}{نمودار تخته: مقدار $PE$ برحسب موقعیت (۰ تا ۱۰۰) برای چهار بعد ۴ تا ۷. هر بعد یک موج با فرکانس ویژهٔ خود است. فرمول $PE(pos,j)=\sin(g(pos,j))$ و یادداشت «هیچ دو خانه‌ای زاویه یکسان ندارند» در تخته بود.}

\begin{definitionbox}{فرمول استاندارد (\eng{Vaswani et al., 2017})}
برای $i=0,1,\dots,d/2-1$:
\[
PE(pos,\,2i)=\sin\!\Big(\frac{pos}{10000^{2i/d}}\Big),\qquad
PE(pos,\,2i+1)=\cos\!\Big(\frac{pos}{10000^{2i/d}}\Big).
\]
یعنی بعدها دو‌به‌دو (سینوس، کسینوس) با یک فرکانس $\omega_i=10000^{-2i/d}$ هستند و با بزرگ‌شدن $i$ فرکانس کم می‌شود: بعدهای کم‌شماره موج‌های سریع (برای تمایز موقعیت‌های نزدیک) و بعدهای پرشماره موج‌های بسیار کند (برای تمایز موقعیت‌های دور) می‌سازند.
\end{definitionbox}

\begin{correctionbox}{فرمول کلاس نسخهٔ ساده‌شدهٔ فرمول استاندارد است}
فرمول $PE(pos,j)=\sin(g(pos,j))$ نسخهٔ ساده‌شدهٔ فرمول استاندارد است: در نسخهٔ کامل \emph{بعدهای فرد کسینوس} می‌گیرند و $g(pos,j)=pos/10000^{2\lfloor j/2\rfloor/d}$. همچنین نمودار تخته برای چهار بعد (۴ تا ۷) دوره‌های حدود ۴۰ و ۱۰۰ را نشان می‌دهد که تصویری آموزشی است؛ در فرمول استاندارد با $d=512$ بعدهای ۴ و ۵ دورهٔ حدود $2\pi/\omega_2\approx 6{,}7$ دارند. کیفیت مشترک هر دو تصویر (جفت‌شدن سینوس و کسینوس و کم‌شدن فرکانس با بعد) درست است. نمودار افزوده‌شدهٔ زیر نمایش واقعی فرمول استاندارد است.

\medskip
دربارهٔ یادداشت «هیچ دو خانه‌ای زاویه یکسان ندارند»: برای هر \emph{بعد} ثابت، دو موقعیت مختلف زاویهٔ متفاوت دارند (چون $\omega_i>0$). اما به‌طور کلی نمی‌توان گفت هیچ دو خانهٔ ماتریس زاویهٔ برابر ندارند؛ مثلاً با $d=512$ داریم $\omega_{i+64}=\omega_i/10$ و بنابراین زاویهٔ ($pos=10$، بعد $i+64$) برابر زاویهٔ ($pos=1$، بعد $i$) است. آنچه مهم است این است که \emph{بردار کامل} هر موقعیت (ترکیب همهٔ بعدها) برای موقعیت‌های مختلف متفاوت است؛ این به‌ویژه به‌خاطر جفت (سینوس، کسینوس) با فرکانس‌های متنوع برقرار است.
\end{correctionbox}

\coursefigure{0.85}{ch10-pe-heatmap.png}{ماتریس $PE$ سینوسی استاندارد با $d=64$ و ۱۰۰ موقعیت (محور افقی: بعد، محور عمودی: موقعیت). بعدهای کم‌شماره (چپ) سریع نوسان می‌کنند و بعدهای پرشماره الگوی کندتر دارند. (شکل افزوده)}

\coursefigure{0.85}{ch10-pe-curves.png}{مقدار $PE$ برحسب موقعیت برای چند بعد از ماتریس $d=64$: فرکانس‌های بالا (ابعاد ۰ و ۱)، متوسط (۸ و ۹) و پایین (۳۲ و ۳۳). هر جفت (سینوس، کسینوس) یک فرکانس دارد. (شکل افزوده)}

\begin{addedbox}{مثال عددی: $d=4$ و موقعیت‌های ۰، ۱، ۲}
برای $d=4$ دو فرکانس داریم: $i=0$ با $\omega_0=1$ و $i=1$ با $\omega_1=10000^{-2/4}=0{,}01$. ردیف هر موقعیت $(\sin(pos\,\omega_0),\cos(pos\,\omega_0),\sin(pos\,\omega_1),\cos(pos\,\omega_1))$ است:
\begin{center}
\begin{tabular}{ccccc}
\toprule
$pos$ & بعد ۰ & بعد ۱ & بعد ۲ & بعد ۳ \\
\midrule
۰ & $0{,}0000$ & $1{,}0000$ & $0{,}0000$ & $1{,}0000$ \\
۱ & $0{,}8415$ & $0{,}5403$ & $0{,}0100$ & $0{,}9999$ \\
۲ & $0{,}9093$ & $-0{,}4161$ & $0{,}0200$ & $0{,}9998$ \\
\bottomrule
\end{tabular}
\end{center}
مثلاً $\sin(1)\approx 0{,}8415$ و $\cos(2)\approx -0{,}4161$. همهٔ مقدارها در $[-1,1]$ و سه ردیف متفاوت‌اند. اگر این ردیف‌ها با \eng{Embedding} ۴ بعدی هر توکن جمع شوند، همان $X=E+P$ به دست می‌آید.
\end{addedbox}

\begin{addedbox}{ویژگی‌های دیگر فرمول سینوسی}
\begin{itemize}
  \item \textbf{قابل محاسبه برای هر موقعیت}: بردار موقعیت ۵۰ بدون مشاهدهٔ آن در دادهٔ آموزشی از فرمول به دست می‌آید.
  \item \textbf{بدون پارامتر آموزشی}: ماتریس با فرمول ساخته می‌شود و گرادیان نمی‌گیرد.
  \item \textbf{ساختار نسبی}: چون $\sin(a+b)$ و $\cos(a+b)$ با فرمول جمع زاویه‌ها نوشته می‌شوند، $PE(pos+k)$ برای هر $k$ ثابت یک تبدیل خطی ثابت (دوران در هر جفت بعد) از $PE(pos)$ است. در نتیجه ضرب داخلی دو بردار موقعیت فقط به فاصلهٔ آن‌ها بستگی دارد: $\frac1d PE(p)\!\cdot\!PE(q)=\frac1d\sum_i\cos(\omega_i(p-q))$.
\end{itemize}
\end{addedbox}

\coursefigure{0.55}{ch10-pe-similarity.png}{شباهت $\frac1d PE(p)\cdot PE(q)$ بین هر دو موقعیت ($d=512$): روی قطر برابر ۱ و با دورشدن از قطر کاهش می‌یابد. مقدار فقط تابع $|p-q|$ است. (شکل افزوده)}

\section{مقایسهٔ دو روش}

\begin{table}[htbp]
\centering
\caption{مقایسهٔ کدگذاری موقعیت یادگرفتنی و ثابت (برگرفته از جدول جلسه، با کمی توضیح افزوده).}
\begin{tabular}{p{3.2cm}p{5.2cm}p{5.2cm}}
\toprule
ویژگی & یادگرفتنی (\eng{Learned}) & ثابت (\eng{Fixed}) \\
\midrule
تعیین مقدارها & از طریق آموزش & با فرمول مشخص \\
پارامتر قابل‌آموزش & دارد ($100\times512$ در مثال) & ندارد \\
وابستگی به دادهٔ آموزشی & زیاد & کم \\
موقعیت‌های دیده‌نشده در آموزش & ممکن است خوب یاد گرفته نشده باشند & با فرمول قابل محاسبه‌اند \\
چالش اصلی & پوشش مناسب موقعیت‌ها در دادهٔ آموزشی & انتخاب قاعدهٔ مناسب برای نمایش موقعیت‌ها \\
انعطاف & زیاد؛ مدل شکل را خودش می‌یابد & محدود به شکل فرمول \\
\bottomrule
\end{tabular}
\end{table}

\section{سؤال مفهومی مهم استاد: دادهٔ آموزشی کوتاه}
\label{sec:ch10-exam}
دو مدل داریم: در مدل اول موقعیت با فرمول محاسبه می‌شود و در مدل دوم موقعیت در طول آموزش یاد گرفته می‌شود. اگر دادهٔ آموزشی عمدتاً جملهٔ کوتاه باشد، برای موقعیت‌های بلندتر کدام انتخاب بهتر است؟

\begin{figure}[htbp]
\centering
\begin{tcolorbox}[colback=softyellow,colframe=black!70,boxrule=0.8pt,arc=6pt,width=0.9\textwidth]
\textbf{دو مدل دارم}
\begin{itemize}
  \item مدل اول: موقعیت با فرمول محاسبه می‌شود.
  \item مدل دوم: موقعیت در طول آموزش یاد گرفته می‌شود.
\end{itemize}
اگر داده‌های آموزشی عمدتاً جمله‌های کوتاه باشند، کدام بهتر است؟
\end{tcolorbox}
\vspace{4pt}
\begin{tcolorbox}[colback=boxgray,colframe=black!40,boxrule=0.5pt,arc=4pt,width=0.9\textwidth]
مدل دوم ممکن است موقعیت‌های پرتکرار را خوب یاد بگیرد، ولی برای موقعیت‌های بسیار دورتر بردار آموزش‌دیده‌ای نداشته باشد. بنابراین، قابل‌آموزش بودن انعطاف بیشتری می‌دهد، اما می‌تواند تعمیم به موقعیت‌های ندیده را دشوار کند.
\end{tcolorbox}
\caption{سؤال (کادر زرد) و پاسخ نوشته‌شده روی تختهٔ کلاس. (بازسازی‌شده از تختهٔ کلاس)}
\end{figure}

\paragraph{پاسخ کلاس.} مدل با \emph{کدگذاری ثابت} انتخاب مناسب‌تری است. دلیل: فرض کنید برای مدل یادگرفتنی ماتریسی با ۱۰۰ سطر داریم ($P\in\mathbb{R}^{100\times 512}$) ولی دادهٔ آموزشی فقط توالی‌هایی با حداکثر طول ۳ دارد. در آموزش فقط سطرهای ۱ تا ۳ استفاده می‌شوند و گرادیان فقط به همین سطرها می‌رسد؛ بردارهای موقعیت ۵۰ یا ۸۰ آموزش مؤثری نمی‌بینند. در آزمون، با یک توالی به طول ۵۰ مدل باید از نمایش‌هایی استفاده کند که هرگز یاد نگرفته است و عملکرد افت می‌کند. در روش ثابت، بردار موقعیت ۵۰ با فرمول محاسبه می‌شود و لازم نیست آن موقعیت در دادهٔ آموزشی دیده شده باشد.

\begin{exambox}{نکات روش‌شناختی این سؤال}
\begin{itemize}
  \item موضوع فقط \emph{تعداد} نمونه‌ها نیست، بلکه \emph{پوشش موقعیت‌ها} در دادهٔ آموزشی است: ممکن است دادهٔ بسیار بزرگ ولی همگی کوتاه باشد، که باز هم تجربهٔ کافی دربارهٔ موقعیت‌های بالاتر نمی‌دهد. اگر دادهٔ آموزشی همهٔ موقعیت‌ها را خوب پوشش دهد، روش یادگرفتنی هم نمایش مناسبی یاد می‌گیرد.
  \item پاسخ دانشجو «اندازهٔ ماتریس را کوچک بگیریم» فرض مسئله را عوض می‌کند. صورت دقیق سؤال این است: \emph{ظرفیت ماتریس ۱۰۰ موقعیت است ولی دادهٔ آموزشی تا طول ۳؛ اندازه را تغییر نمی‌دهیم و در همین شرایط دو روش را مقایسه می‌کنیم.}
  \item استاد گفت در امتحان فرض‌های سؤال را دقیق‌تر بیان خواهد کرد تا چنین ابهامی پیش نیاید.
\end{itemize}
\textbf{صورت پیشنهادی امتحان:} «دو مدل داریم که هر دو برای توالی‌هایی با حداکثر طول ۱۰۰ طراحی شده‌اند. مدل اول از کدگذاری موقعیت ثابت و مدل دوم از کدگذاری یادگرفتنی استفاده می‌کند. اگر تمام دادهٔ آموزشی فقط شامل توالی‌هایی با طول حداکثر ۳ باشد، کدام روش برای مواجهه با موقعیت‌های بلندتر مناسب‌تر است و چرا؟»
\end{exambox}

\begin{correctionbox}{دقت: «ثابت بهتر است» تضمین کامل نیست}
مزیت روش ثابت این است که \emph{خودِ نمایش موقعیت} به آموزش هر سطر وابسته نیست. این به‌تنهایی تضمین نمی‌کند که کل مدل روی توالی‌های بلندتر از آنچه دیده خوب کار کند؛ وزن‌های \eng{Attention} هم فقط روی موقعیت‌های کوتاه آموزش دیده‌اند و در عمل \eng{Transformer}های سینوسی هم در بلندتر از طول آموزش اغلب افت می‌کنند (به همین دلیل روش‌های موقعیت نسبی مطرح شدند). پس پاسخ درست برای امتحان این است که روش ثابت \emph{مناسب‌تر} است، نه اینکه مشکل را کاملاً حل کند. استاد هم همین‌گونه نتیجه گرفت.
\end{correctionbox}

\section{کدگذاری موقعیت به توالی وابسته است، نه به «جمله»}
دانشجو پرسید مدل‌هایی مانند \eng{Transformer} در سطح توکن کار می‌کنند و لزوماً جمله را واحد اصلی نمی‌گیرند. پاسخ استاد: کدگذاری موقعیت ذاتاً به مفهوم «جمله» وابسته نیست. عنصرهای توالی می‌توانند کلمه، نویسه، واج یا حتی بایت باشند و کدگذاری موقعیت فقط مشخص می‌کند هر عنصر \emph{چندمین عنصر توالی ورودی} است:
\begin{itemize}
  \item ورودی نویسه‌ای: موقعیت پنجم یعنی نویسهٔ پنجم؛
  \item ورودی کلمه‌ای: موقعیت پنجم یعنی کلمهٔ پنجم؛
  \item ورودی واجی: موقعیت پنجم یعنی واج پنجم.
\end{itemize}
اگر ورودی شبکه ۱۰۰ نویسه باشد، ۱۰۰ موقعیت به ۱۰۰ نویسه برمی‌گردد؛ اگر ۱۰۰ کلمه باشد، همان تعداد موقعیت برای کلمه‌ها به کار می‌رود. این کاملاً به طراحی معماری و \eng{Tokenization} بستگی دارد.

\section{طول زمینه (\eng{Context Length})}
پرسش دانشجو: اگر توالی ۱۱۰ عنصر باشد ولی ماتریس موقعیت فقط برای ۱۰۰ موقعیت تعریف شده باشد چه می‌شود؟

\begin{itemize}
  \item حداکثر طول ورودی یک تصمیم طراحی است. برای توالی‌های بلندتر باید تصمیم بگیریم چه کنیم؛ یک راه شکستن ورودی به بخش‌های کوچک‌تر و استفاده از پنجره (\eng{Window} یا \eng{Context Window}) است، یعنی مدل در هر گام بخش مشخصی از توالی را پردازش می‌کند.
  \item اندازهٔ پوشش موقعیت به طراحی مدل و \eng{Task} بستگی دارد و یک عدد ثابت برای همهٔ مدل‌ها و مسئله‌ها وجود ندارد. مثلاً مدل نویسه‌ای معمولاً به موقعیت‌های بیشتری نیاز دارد تا مدل کلمه‌ای، چون توالی نویسه‌ها بلندتر است.
  \item برای ورودی‌های چند هزار یا چند ده هزار عنصری، علاوه بر پنجره‌بندی، معماری‌های مختلف ماژول‌ها و شیوه‌های متفاوتی برای مدیریت اطلاعات موقعیت دارند؛ محل و تعداد این ماژول‌ها بخشی از طراحی معماری است.
\end{itemize}

\begin{addedbox}{دو رفتار متفاوت برای عنصر ۱۰۱ام}
در روش \emph{یادگرفتنی} با جدول ۱۰۰ سطری، موقعیت ۱۰۱ سطری ندارد و عملاً خطای اندیس رخ می‌دهد؛ ورودی باید بریده یا تقسیم شود. در روش \emph{فرمولی} می‌توان $PE(101)$ را حساب کرد، ولی مدل با چنین موقعیتی هرگز آموزش ندیده و کیفیت تضمین‌شده نیست. این مثال دوباره نشان می‌دهد چرا طول حداکثر یک فراپارامتر معماری است. روش‌های جدیدتر (کدگذاری موقعیت نسبی، \eng{RoPE}، \eng{ALiBi}) برای کاهش این محدودیت پیشنهاد شده‌اند؛ خارج از محتوای کلاس.
\end{addedbox}

\section{مثال عددی کامل: \eng{Attention} در رمزگشا}
\label{sec:ch10-example}
در ادامهٔ جلسه، استاد روی تخته یک مثال عددی دست‌نویس از محاسبات \eng{Attention} در حالت رمزگذار--رمزگشا (همان معماری \eng{G2P}) حل کرد. تصویر بالای تخته دو پشته را نشان می‌دهد: رمزگذار (ورودی نویسه‌ها با \eng{Embedding} نویسه، \eng{CE}، و کدگذاری موقعیت \eng{PE}) و رمزگشا (توکن فعلی با \eng{CE} و \eng{PE}) که \eng{Query} آن بردار بافت $C$ را می‌سازد؛ بالای تخته برچسب چند واج نوشته شده است.

\begin{figure}[htbp]
\centering
\resizebox{0.95\textwidth}{!}{%
\begin{tikzpicture}[rtl, 
  cell/.style={draw,rounded corners=2pt,minimum width=1.3cm,minimum height=0.7cm,font=\small},
  blk/.style={draw,rounded corners=4pt,fill=orange!12,align=center,font=\normalsize},
  ar/.style={-{Stealth[length=2.2mm]},thick}]
  % encoder
  \node[font=\large\bfseries] at (2.0,-2.3) {\RL{رمزگذار}};
  \foreach \i/\x in {1/0,2/2,3/4}{
    \node (x\i) at (\x,-1.4) {$x_{\i}$};
    \node[cell,fill=softblue] (ce\i) at (\x,0) {\lr{CE}};
    \node[cell,fill=softgreen] (pe\i) at (\x,1.2) {\lr{PE}};
    \node at (\x,0.6) {\small $+$};
    \draw[ar] (x\i)--(ce\i);
    \draw[ar] (pe\i.north)--(\x,2.1);
  }
  \node[blk,minimum width=5.6cm,minimum height=1.2cm] (enc) at (2,2.7) {\lr{Encoder}};
  % decoder
  \node[font=\large\bfseries] at (10.5,-2.3) {\RL{رمزگشا}};
  \node (y) at (10.5,-1.4) {\lr{$y_{t-1}$}};
  \node[cell,fill=softblue] (dce) at (10.5,0) {\lr{CE}};
  \node[cell,fill=softgreen] (dpe) at (10.5,1.2) {\lr{PE}};
  \node at (10.5,0.6) {\small $+$};
  \draw[ar] (y)--(dce);
  \node[blk,minimum width=3.6cm,minimum height=1.2cm,fill=orange!20] (att) at (10.5,3.4) {\lr{Attention}};
  \draw[ar] (dpe.north)--node[right,font=\small]{$Q$} (att.south);
  \draw[ar] (enc.east)--node[above,font=\small]{$K,\;V$} (att.west |- enc.east);
  \node[draw,rounded corners=3pt,fill=softyellow,minimum width=3cm] (C) at (10.5,5.3) {$C$\ \lr{Context Vector}};
  \draw[ar] (att)--(C);
  \node[draw,rounded corners=3pt,fill=softgreen,minimum width=3cm] (out) at (10.5,6.8) {\lr{Logits + Softmax}};
  \draw[ar] (C)--(out);
  \node[font=\normalsize] (ph) at (10.5,8.1) {\RL{برچسب واج خروجی}};
  \draw[ar] (out)--(ph);
\end{tikzpicture}}
\caption{نمای کلی مثال تخته: ورودی رمزگذار (\eng{CE} به‌علاوهٔ \eng{PE}) و ورودی رمزگشا (\eng{CE} به‌علاوهٔ \eng{PE}) که \eng{Query} آن با \eng{Key}/\eng{Value} رمزگذار ترکیب می‌شود و بردار بافتی $C$ را می‌سازد. برچسب‌های واجِ بالای تخته دست‌خط ناخوانا بود و اینجا فقط به‌صورت کلی نشان داده شده‌اند. (بازسازی‌شده از تختهٔ کلاس)}
\end{figure}

\subsection{گام ۱: \eng{Key} و \eng{Value} از خروجی رمزگذار}
سه توکن ورودی داریم و ماتریس \eng{Embedding} آن‌ها $E\in\mathbb{R}^{3\times 2}$ است. با ضرب در $W_K$ و $W_V$ (هر دو $2\times2$) ماتریس‌های $K$ و $V$ هر دو $3\times2$ به دست می‌آیند. مقدارهای نوشته‌شده روی تخته در شکل زیر آمده است:

\begin{figure}[htbp]
\centering
\resizebox{0.95\textwidth}{!}{%
\begin{tikzpicture}[rtl, ar/.style={-{Stealth[length=2.5mm]},thick,mainblue},
  fr/.style={draw,minimum width=1.2cm,minimum height=1.1cm}]
  \node (E) at (0,0) {$\underset{3\times2}{E=\begin{bmatrix}1&2\\3&4\\5&6\end{bmatrix}}$};
  \node (K) at (9,1.9) {$\underset{3\times2}{K=\begin{bmatrix}9&10\\11&12\\13&14\end{bmatrix}}$};
  \node (V) at (9,-1.9) {$\underset{3\times2}{V=\begin{bmatrix}15&16\\17&18\\19&20\end{bmatrix}}$};
  \draw[ar] (1.9,0.5)--node[above,sloped]{$\times\,\underset{2\times2}{W_K}$} (6.6,1.9);
  \draw[ar] (1.9,-0.5)--node[above,sloped]{$\times\,\underset{2\times2}{W_V}$} (6.6,-1.9);
  \node[draw=red!70!black,circle,thick,inner sep=2pt] at (-2.4,0) {۱};
\end{tikzpicture}}
\caption{گام ۱: $E$ ($3\times2$) در $W_K$ و $W_V$ ($2\times2$) ضرب می‌شود و $K$ و $V$ ($3\times2$) به دست می‌آیند. روی تخته $W_K$ و $W_V$ فقط به‌صورت قاب‌های خالی کشیده شده بودند و مقدارشان نوشته نشده بود. (بازسازی‌شده از تختهٔ کلاس)}
\end{figure}

\begin{addedbox}{ماتریس‌های وزنی سازگار با $K$ و $V$}
چون سطرهای $E$ تصاعد حسابی‌اند، ماتریس‌های $2\times2$ یکتا و سازگاری وجود دارند: حل دستگاه $EW_K=K$ از دو سطر اول و بررسی سطر سوم ($5\cdot(-7)+6\cdot 8=13$ و $5\cdot(-8)+6\cdot 9=14$) می‌دهد
\[
W_K=\begin{bmatrix}-7&-8\\8&9\end{bmatrix},\qquad
W_V=\begin{bmatrix}-13&-14\\14&15\end{bmatrix}.
\]
یعنی $K=EW_K$ و $V=EW_V$ دقیقاً برقرارند.
\end{addedbox}

\subsection{گام ۲: \eng{Query} از توکن رمزگشا}
توکن فعلی رمزگشا بعد از \eng{Embedding} با کدگذاری موقعیتش جمع می‌شود: $x=E_{dec}+P_{pos}$ (برداری $1\times2$). سپس با $W_Q$ ($2\times2$) به \eng{Query} ($1\times 2$) تبدیل می‌شود:
\[
Q = (E_{dec}+P_{pos})\,W_Q\in\mathbb{R}^{1\times 2}.
\]
یادداشت زردرنگ روی تخته در همین جا «$+\,PE$» را نشان می‌دهد؛ یعنی موقعیت قبل از ساخت $Q$ در نمایش وارد می‌شود.

\subsection{گام ۳: امتیازها $Q K^{\top}$}
\[
\text{\lr{scores}} = Q K^{\top}\in\mathbb{R}^{1\times 3}\quad(1\times2)\cdot(2\times3).
\]
\begin{figure}[htbp]
\centering
\resizebox{0.95\textwidth}{!}{%
\begin{tikzpicture}[rtl]
  \node[draw=red!70!black,circle,thick,inner sep=2pt] at (-0.5,0) {۲};
  \node[draw=orange,rounded corners=3pt,fill=orange!12,inner sep=5pt] (ep) at (2.5,0) {$\underset{1\times2}{(E_{dec}+PE)}$};
  \node at (5.1,0) {$\times$};
  \node[draw,minimum width=1.3cm,minimum height=1.1cm,label=below:{\small \lr{$W_Q$ ($2\times2$)}}] at (6.5,0) {};
  \node at (8.0,0) {$=$};
  \node (Q) at (10.3,0) {$\underset{1\times2}{Q=\begin{bmatrix}1&2\end{bmatrix}}$};
  \node[draw=red!70!black,circle,thick,inner sep=2pt] at (-0.5,-3) {۳};
  \node (sc) at (7.2,-3) {$\underset{1\times2}{Q}\;\cdot\;\underset{2\times3}{K^{\top}=\begin{bmatrix}9&11&13\\10&12&14\end{bmatrix}}\;=\;\underset{1\times3}{\begin{bmatrix}29&35&41\end{bmatrix}}$};
  \node[font=\small,text=gray] at (7.2,-4.3) {\RL{امتیازها} \lr{scores}};
\end{tikzpicture}}
\caption{گام‌های ۲ و ۳: \eng{Embedding} توکن رمزگشا با \eng{PE} جمع و در $W_Q$ ضرب می‌شود تا $Q$ ساخته شود (بالا)؛ سپس $QK^{\top}$ امتیازها را می‌دهد (پایین). مقدارهای $Q$ و امتیازها دست‌خط به‌سختی خوانایی بودند و در اینجا بر پایهٔ خوانش توضیح‌داده‌شده در متن آمده‌اند. (بازسازی‌شده از تختهٔ کلاس)}
\end{figure}

\begin{correctionbox}{نکتهٔ خوانایی: مقدارهای $Q$ و امتیازها}
اعداد دست‌نویس $Q$ و امتیازها در تخته به‌سختی خوانا هستند. با توجه به ساختار $K^\top=\begin{bmatrix}9&11&13\\10&12&14\end{bmatrix}$، خوانش‌ای که با اعداد تخته بیشترین سازگاری را دارد $Q=(1,2)$ و امتیازها $(29,35,41)$ است:
\[
QK^\top=(1\cdot9+2\cdot10,\;1\cdot11+2\cdot12,\;1\cdot13+2\cdot14)=(29,\,35,\,41).
\]
این خوانش قطعی نیست؛ ولی ماتریس‌های $E$، $K$ و $V$ و ابعاد همهٔ گام‌ها روی تخته کاملاً خوانا بودند و منطق محاسبه از مقدار دقیق $Q$ مستقل است. در ادامه برای تکمیل محاسبه از همین $Q=(1,2)$ استفاده می‌شود.
\end{correctionbox}

\subsection{گام ۴: بردار بافت (جمع وزن‌دار \eng{Value}ها)}
\[
C=\sum_{i=1}^{3}\alpha_i V_i=\alpha_1(15,16)+\alpha_2(17,18)+\alpha_3(19,20)\in\mathbb{R}^{1\times 2}.
\]
روی تخته سه ضریب کنار سه سطر $V$ نوشته شد و بردار حاصل با برچسب \eng{Context Vector} دایره شد. ضریب‌های نوشته‌شده روی تخته با خوانش بالا به امتیازها شبیه‌اند و زیر آن‌ها یک یادداشت آبی به‌صورت جمع سه عدد دیده می‌شود؛ ولی خوانش دقیق این اعداد ممکن نیست، پس در ادامه صرفاً روش درست محاسبه توضیح داده می‌شود.

\begin{correctionbox}{امتیاز خام ضریب توجه نیست}
در \eng{Attention} واقعی، امتیازهای $QK^\top$ پیش از ضرب در $V$ باید به وزن‌هایی با مجموع ۱ تبدیل شوند (\eng{Softmax}، معمولاً پس از تقسیم بر $\sqrt{d_k}$). اگر امتیازهای خام ضریب شوند، حاصل
\[C=29(15,16)+35(17,18)+41(19,20)=(1809,\,1914)\]
می‌شود که از دامنهٔ مقدارهای $V$ ($15$ تا ۲۰) بسیار بزرگ‌تر است. حتی تقسیم بر مجموع ($29+35+41=105$) نتیجهٔ $C\approx(17{,}23;\,18{,}23)$ را می‌دهد که قابل قبول‌تر است ولی هنوز \eng{Softmax} نیست.
\end{correctionbox}

\begin{addedbox}{محاسبهٔ درست با \eng{Softmax}}
با امتیازهای $(29,35,41)$:
\[
\begin{aligned}
\alpha&=\mathrm{softmax}(29,35,41)\approx(0{,}000006;\;0{,}0025;\;0{,}9975),\\
C&=\alpha V\approx(18{,}995;\;19{,}995).
\end{aligned}
\]
چون اختلاف امتیازها بزرگ است، \eng{Softmax} تقریباً همهٔ وزن را به یک عنصر می‌دهد (توجه «تیز»). برای همین امتیازها را بر $\sqrt{d_k}=\sqrt2$ تقسیم می‌کنند:
$\mathrm{softmax}(20{,}51;\,24{,}75;\,28{,}99)\approx(0{,}0002;\;0{,}0142;\;0{,}9856)$ و $C\approx(18{,}97;\;19{,}97)$. در یک مدل واقعی با $d_k$ بزرگ‌تر این مقیاس‌دهی مهم‌تر است.
\end{addedbox}

\subsection{گام ۵: لجیت و \eng{Softmax} روی واژگان}
بردار بافت $C\,(1\times2)$ از یک لایهٔ خروجی می‌گذرد:
\[
\text{\lr{logits}} = C\,W + b,\qquad W\in\mathbb{R}^{2\times 5},\; b\in\mathbb{R}^{1\times 5},
\]
که لجیت $1\times 5$ می‌دهد (روی تخته اندازه‌ها $2\times5$، $1\times5$ و $1\times5$ نوشته شدند؛ یعنی اندازهٔ واژگانِ مثال ۵ است). در پایان \eng{Softmax} روی واژگان احتمال هر توکن بعدی را می‌دهد.

\begin{figure}[htbp]
\centering
\resizebox{0.95\textwidth}{!}{%
\begin{tikzpicture}[rtl]
  \node[draw=red!70!black,circle,thick,inner sep=2pt] at (-0.5,0) {۴};
  \node (C) at (8,0) {$C=\alpha_1\begin{bmatrix}15&16\end{bmatrix}+\alpha_2\begin{bmatrix}17&18\end{bmatrix}+\alpha_3\begin{bmatrix}19&20\end{bmatrix}\;=\;\underset{1\times2}{\boxed{\;C\;}}$};
  \node[font=\small,text=gray] at (14.0,-1.1) {\lr{Context Vector}};
  \node[draw=red!70!black,circle,thick,inner sep=2pt] at (-0.5,-3.4) {۵};
  \node (L) at (7.2,-3.4) {$\underset{1\times5}{\text{\lr{logits}}}\;=\;\underset{1\times2}{C}\;\underset{2\times5}{W}\;+\;\underset{1\times5}{b}\;=\;\begin{bmatrix}z_1&z_2&z_3&z_4&z_5\end{bmatrix}\;\xrightarrow{\;\text{\lr{softmax}}\;}\;\underset{1\times5}{\begin{bmatrix}p_1&p_2&p_3&p_4&p_5\end{bmatrix}}$};
\end{tikzpicture}}
\caption{گام‌های ۴ و ۵: جمع وزن‌دار سطرهای $V$ با ضرایب $\alpha_i$ بردار بافت $C$ ($1\times2$) را می‌سازد؛ سپس لجیت $CW+b$ ($1\times5$) و \eng{Softmax} به احتمال هر توکن می‌رسد. ضرایب نوشته‌شده روی تخته ناخوانا بود و به‌صورت نماد $\alpha_i$ آمده است. (بازسازی‌شده از تختهٔ کلاس)}
\end{figure}

\begin{addedbox}{نمونهٔ عددی گام ۵}
فرض کنید $C\approx(19,20)$ و
$W=\begin{bmatrix}0{,}1&-0{,}05&0{,}02&0{,}04&-0{,}1\\0&0{,}05&-0{,}03&0{,}06&0{,}1\end{bmatrix}$ و $b=0$. آنگاه
$\text{\lr{logits}}=(1{,}90;\,0{,}05;\,-0{,}22;\,1{,}96;\,0{,}10)$ و
$\mathrm{softmax}\approx(0{,}399;\,0{,}063;\,0{,}048;\,0{,}424;\,0{,}066)$ که مجموعش ۱ است و توکن چهارم محتمل‌ترین است. (مقدارهای $W$ فرضی‌اند.)
\end{addedbox}

\subsection{ماسک علّی (پایین‌مثلثی)}
در سمت راست تخته، $QK^{\top}$ برای همهٔ موقعیت‌های رمزگشا (ماتریس مربعی با برچسب سطرها و ستون‌ها) و یک ماتریس ماسک نیز کشیده شد: مثلث بالا‌راست خط‌خورده (\emph{ماسک‌شده}) است و ناحیهٔ پایین‌چپ باقی می‌ماند؛ طبق روش معمول، قطر اصلی هم باقی می‌ماند (هر توکن به خودش و توکن‌های قبلی نگاه می‌کند).

\begin{figure}[htbp]
\centering
\begin{tikzpicture}[rtl, x=1.3cm,y=1.0cm]
  \foreach \i in {1,...,4}{
    \node at (\i-0.5,0.5) {$t_{\i}$};
    \node[left] at (-0.1,-\i+0.5) {$t_{\i}$};
    \foreach \j in {1,...,4}{
      \ifnum\j>\i
        \fill[softred] (\j-1,-\i+1) rectangle (\j,-\i);
        \draw[red!60!black,thin] (\j-1,-\i) -- (\j,-\i+1);
        \draw[red!60!black,thin] (\j-1,-\i+0.5) -- (\j-0.5,-\i+1);
        \draw[red!60!black,thin] (\j-0.5,-\i) -- (\j,-\i+0.5);
        \node[font=\footnotesize] at (\j-0.5,-\i+0.5) {$-\infty$};
      \else
        \fill[softgreen] (\j-1,-\i+1) rectangle (\j,-\i);
        \node at (\j-0.5,-\i+0.5) {$s_{\i\j}$};
      \fi
      \draw[gray] (\j-1,-\i+1) rectangle (\j,-\i);
    }
  }
  \node[above=0.5cm,font=\small] at (2,1) {\lr{Key}};
  \node[font=\small,rotate=90] at (-1.3,-2) {\lr{Query}};
  \node[right,font=\small,align=left] at (4.4,-1) {\RL{ناحیهٔ خط‌خورده: ماسک‌شده}\\\RL{(توکن‌های آینده)}};
  \node[right,font=\small,align=left] at (4.4,-3) {\RL{ناحیهٔ سبز: مجاز}\\\RL{(قطر و توکن‌های قبلی)}};
\end{tikzpicture}
\caption{ماتریس امتیازهای $QK^{\top}$ همهٔ موقعیت‌های رمزگشا و ماسک علّی: مثلث بالا‌راست (خط‌خورده) پیش از \eng{Softmax} با $-\infty$ جایگزین می‌شود. روی تخته چهار موقعیت کشیده شده بود و برچسب‌های آن ناخوانا بود. (بازسازی‌شده از تختهٔ کلاس)}
\end{figure}

دلیل: هنگام تولید توکن موقعیت $t$ رمزگشا نباید توکن‌های آینده را ببیند، چون هنگام تولید واقعی هنوز وجود ندارند. با افزودن $-\infty$ به امتیازهای آینده پیش از \eng{Softmax}، وزن آن‌ها صفر می‌شود. برای ۳ توکن:
\[
\text{\lr{mask}}=\begin{bmatrix}0&-\infty&-\infty\\0&0&-\infty\\0&0&0\end{bmatrix},\qquad
\alpha=\mathrm{softmax}\big(QK^\top/\sqrt{d_k}+\text{\lr{mask}}\big).
\]
\begin{exambox}{یادداشت‌های پایانی روی تخته}
جدا از مثال، روی چارچوب \eng{TTS} دو یادداشت جدید آمده بود: (۱) همان پرسش بحث دربارهٔ نیاز روش بازگشتی به کدگذاری موقعیت (بخش \ref{sec:ch10-rnn}) و (۲) «مدل‌های اتورگرسیو (\eng{Autoregressive}) و \eng{Teacher Forcing} در زمان آموزش». مورد دوم هنوز در این جلسه تدریس نشد و موضوع بعدی است.
\end{exambox}

\begin{addedbox}{مدل اتورگرسیو و \eng{Teacher Forcing} (پیش‌نمایش)}
در مدل \term{اتورگرسیو}{Autoregressive} هر خروجی به خروجی‌های پیشین وابسته است: $P(y_1,\dots,y_T)=\prod_t P(y_t\mid y_{<t},x)$. هنگام استنتاج مدل خروجی قبلی خودش را به‌عنوان ورودی بعدی می‌دهد. هنگام آموزش با \eng{Teacher Forcing} به‌جای خروجی پیش‌بینی‌شدهٔ مدل، توکن \emph{درست} موقعیت قبلی به رمزگشا داده می‌شود. این باعث می‌شود همهٔ موقعیت‌ها هم‌زمان (موازی) با کمک ماسک علّی آموزش ببینند و خطای یک گام روی گام‌های بعد پخش نشود؛ اما شکاف آموزش و استنتاج هم ایجاد می‌کند. جزئیات در ادامهٔ درس خواهد آمد.
\end{addedbox}

\section{جمع‌بندی}
\begin{itemize}
  \item \eng{Self-Attention} موازی بدون اطلاعات موقعیت نسبت به ترتیب عنصرها بی‌اطلاع است (هم‌متغیر نسبت به جایگشت)؛ مدل \eng{RNN} ترتیب را ضمناً از بازگشت می‌گیرد.
  \item کدگذاری موقعیت پس از \eng{Embedding} اولیه و پیش از \eng{Attention} با جمع $X=E+P$ اعمال می‌شود؛ جمع، ابعاد را ثابت نگه می‌دارد.
  \item ماتریس $P\in\mathbb{R}^{L\times d}$ را می‌توان یاد گرفت (جدول جست‌وجو، بدون فعال‌ساز یا ضرب ماتریسی، فقط سطرهای دیده‌شده گرادیان می‌گیرند) یا با فرمول سینوسی ساخت (بدون پارامتر، برای هر موقعیت قابل محاسبه).
  \item موقعیت به توالی ورودی (کلمه، نویسه، واج، بایت) مربوط است نه لزوماً «جمله»؛ حداکثر طول توالی یک تصمیم طراحی است.
  \item در مثال عددی، $K$ و $V$ از رمزگذار، $Q$ از رمزگشا (با $+PE$)، امتیازها $QK^\top$، \eng{Softmax}، $C=\alpha V$، لجیت $CW+b$ و ماسک علّی به‌ترتیب آمدند.
\end{itemize}

\begin{exambox}{نکات مهم برای مرور و امتحان}
\begin{enumerate}
  \item چرا \eng{Self-Attention} بدون کدگذاری موقعیت دو جملهٔ «ایران آمریکا را شکست داد» و «آمریکا ایران را شکست داد» را با هم اشتباه می‌گیرد؟ چرا \eng{RNN} این مشکل را ندارد؟
  \item محل کدگذاری موقعیت در معماری، و دلیل آن (پیش از \eng{Attention}).
  \item مقایسهٔ کانکت و جمع؛ فرمول $X=E+P$ و اندازهٔ $P\in\mathbb{R}^{100\times512}$؛ تفسیر سطر و ستون.
  \item کدگذاری یادگرفتنی: پارامتر قابل‌آموزش، یادگیری با \eng{Backprop}، بدون فعال‌ساز و ضرب ماتریسی (لایهٔ خطی نیست).
  \item کدگذاری ثابت: ورودی‌های فرمول ($pos$ و $j$)، کران $[-1,1]$ سینوس و کسینوس.
  \item سؤال امتحانی مهم: دادهٔ آموزشی کوتاه (طول ۳) و ماتریس ۱۰۰ موقعیتی؛ چرا کدگذاری ثابت مناسب‌تر است؟ در پاسخ فرض‌ها را (ظرفیت ۱۰۰، دادهٔ تا طول ۳، بدون تغییر اندازهٔ ماتریس) دقیق بنویسید.
  \item ممکن است سؤال طراحی بیاید: اطلاعات موقعیت را در کجای یک شبکهٔ داده‌شده وارد می‌کنید؟
  \item مثال عددی \eng{Attention}: ابعاد هر گام ($E,K,V:3\times2$، $Q:1\times2$، $QK^\top:1\times3$، $C:1\times2$، لجیت $1\times5$).
\end{enumerate}
\end{exambox}
